problem:  
Let \((M,\mathcal{D},g)\) be a compact, smooth, **equiregular step‑3 sub‑Riemannian manifold** equipped with its Popp measure, and let \(\Delta\) denote the intrinsic (Popp‑symmetric) sub‑Laplacian with heat kernel \(p_t(x,y)\). For step‑2 examples such as contact or quasi‑contact manifolds, the on‑diagonal heat kernel admits a full small‑time expansion  
\[
p_t(x,x)\sim t^{-Q/2}\sum_{j=0}^\infty a_j(x)t^j,\qquad t\to0^+,
\]
proved using nilpotent approximation and the Heisenberg/filtered pseudodifferential calculi.  

**Open problem:**  
Determine the precise analytic structure of the small‑time diagonal asymptotic of \(p_t(x,y)\) on equiregular step‑3 sub‑Riemannian manifolds. In particular:  

1. Is the asymptotic necessarily a **pure power‑series expansion in integer powers of \(t\)**, as in step‑2 geometries, or can non‑integer powers or logarithmic terms appear?  
2. Equivalently, is the heat kernel on an equiregular step‑3 sub‑Riemannian manifold modelled by its tangent Carnot group up to all orders in a smooth asymptotic, or does failure of homogeneous dilation symmetry produce analytic anomalies (e.g., \(t^{m/q}\) or \(t^{m}\log t\) terms)?  
3. Can one construct (or rule out) explicit examples where such anomalous terms arise, perhaps in structures whose nilpotent model has rank growth \(2\!\to\!3\!\to\!n\) with non‑trivial brackets of length three?  

Why is it a "good" problem:  
This problem asks for a structural classification of heat kernel asymptotics beyond the well‑understood step‑2 setting. Step‑3 manifolds lie just past the threshold where existing geometric calculi (Heisenberg, Rockland, or filtered \(Ψ\)‑DO frameworks) guarantee smooth power‑series behavior, yet no universal theorem rules out the appearance of non‑integer or logarithmic corrections. Resolving this would clarify whether the **filtered pseudodifferential approach** developed by van Erp–Yuncken and successors fully controls all equiregular geometries or whether additional analytic phenomena emerge in higher step.  

A proof that the expansion is always polynomial would extend the reach of the current hypoelliptic calculus and unify sub‑Riemannian microlocal analysis; a counterexample would reveal fundamentally new singular behavior of the heat kernel and prompt a refinement of existing symbolic calculi.  

Its formulation is simple—“Does the step‑3 heat kernel still have a power‑series expansion?”—yet its resolution demands deep integration of **geometric analysis**, **filtered differential operators**, and **nilpotent approximation theory**, epitomizing a “narrow entrance, profound depth” research problem.